\chapter{Guida alla lettura}

\section*{Organizzazione}

Gli appunti seguono l'ordine di studio, non quello delle lezioni. Ogni concetto è definito una sola volta, nel
punto in cui serve per la prima volta. Dopo viene solo richiamato con un rimando.

\begin{itemize}
  \item \textbf{Parte I --- Prodotti software e sviluppo agile.} Che cosa cambia quando si sviluppa un prodotto
    invece di un progetto (capitolo~\ref{cap:prodotti}). Come lo si sviluppa con i metodi agili
    (capitolo~\ref{cap:agile}). Come si decide che cosa fa (capitolo~\ref{cap:feature}).
  \item \textbf{Parte II --- Architettura del software.} Come si organizza un prodotto in componenti, strati e
    server, e quali scelte tecnologiche lo vincolano (capitolo~\ref{cap:architettura}).
  \item \textbf{Parte III --- Cloud e microservizi.} I container e Docker (capitolo~\ref{cap:container}), il software
    che gira sul cloud (capitolo~\ref{cap:cloud}) e l'architettura a microservizi (capitolo~\ref{cap:microservizi}).
    Infine gli smell architetturali dei microservizi e i refactoring che li risolvono (capitolo~\ref{cap:smell}).
  \item \textbf{Parte IV --- Dal monolite alla produzione.} Come si scompone un monolite in microservizi, database
    compreso (capitolo~\ref{cap:monolite}). Kubernetes, la piattaforma che orchestra i container in produzione
    (capitolo~\ref{cap:kubernetes}). Come si collauda un prodotto (capitolo~\ref{cap:testing}). Come un unico team lo
    sviluppa, lo rilascia e lo gestisce con l'automazione (capitolo~\ref{cap:devops}). Come lo si protegge dagli attacchi
    e come si rispetta la privacy (capitolo~\ref{cap:sicurezza}). Poi i problemi di sicurezza specifici dei
    microservizi (capitolo~\ref{cap:sicurezza-micro}). Infine il quantum computing e le tecniche di Quantum Software
    Engineering per eseguire meglio i programmi quantistici (capitolo~\ref{cap:quantum}).
\end{itemize}

Ogni capitolo si apre con un breve riquadro che ne descrive il contenuto e si chiude con i \emph{punti chiave}
e alcune \emph{domande di verifica}.

Accanto al titolo di alcune sezioni, esempi e tabelle compare la nota piccola \aggiunta. Indica un approfondimento
che spiega meglio un concetto, con esempi e passaggi in più.

\section*{Convenzioni grafiche}

\begin{definizione}[nome]
Definizioni dei concetti.
\end{definizione}
\begin{idea}
Il concetto che bisogna aver capito di una sezione.
\end{idea}
\begin{attenzione}
Errori frequenti e affermazioni da non fraintendere. È l'unico riquadro colorato.
\end{attenzione}
\begin{esempio}
Esempi svolti e casi concreti.
\end{esempio}

Nel testo il \textbf{grassetto} segnala un termine definito per la prima volta. Il \ingl{corsivo} segnala i termini
in un'altra lingua e le enfasi. Il \uline{sottolineato} segnala le poche affermazioni da ricordare alla lettera.
I termini tecnici non sono tradotti: restano in inglese.

\section*{Sigle e simboli}

\begin{center}
\small
\begin{tabularx}{0.92\linewidth}{>{\raggedright\arraybackslash}p{3.2cm}X}
\toprule
Sigla & Significato \\
\midrule
PBI & \ingl{product backlog item}, voce del product backlog (\rif{sec:agile-backlog}) \\
XP & \ingl{Extreme Programming} (\rif{sec:agile-xp}) \\
SOA & \ingl{service-oriented architecture}, architettura orientata ai servizi (\rif{sec:arch-distribuzione}) \\
MVC & pattern \ingl{Model-View-Controller} (\rif{sec:arch-distribuzione}) \\
IaaS, PaaS, SaaS, FaaS & \ingl{Infrastructure, Platform, Software, Function as a Service} (\rif{sec:cloud-aas}) \\
VM & \ingl{virtual machine}, macchina virtuale \\
API & \ingl{application programming interface}, interfaccia con cui altri componenti usano i servizi \\
$C$, $A$, $P$ & \ingl{Consistency}, \ingl{Availability}, \ingl{Partition-tolerance} (\rif{sec:micro-dati}) \\
K8s & Kubernetes (capitolo~\ref{cap:kubernetes}) \\
TDD & \ingl{test-driven development} (\rif{sec:test-tdd}) \\
CI, CD & \ingl{continuous integration}, \ingl{continuous delivery/deployment} (\rif{sec:devops-ci}, \rif{sec:devops-cd}) \\
IaC, CM & \ingl{infrastructure as code}, \ingl{configuration management} (\rif{sec:devops-iac}) \\
APM & \ingl{application performance monitoring} (\rif{sec:devops-approccio}) \\
DoS, DDoS & \ingl{denial of service}, \ingl{distributed denial of service} (\rif{sec:sic-dos}) \\
XSS & \ingl{cross-site scripting} (\rif{sec:sic-xss}) \\
ACL & \ingl{access control list} (\rif{sec:sic-autorizzazione}) \\
TLS, CA & \ingl{Transport Layer Security}, \ingl{certificate authority} (\rif{sec:sic-tls}) \\
KMS & \ingl{key management system} (\rif{sec:sic-chiavi}) \\
JWT & \ingl{JSON Web Token} (\rif{sec:secm-sfide}) \\
QPU, QSE & \ingl{Quantum Processing Unit}, \ingl{Quantum Software Engineering} (\rif{sec:qse-qse}) \\
NISQ & \ingl{Noisy Intermediate-Scale Quantum} (\rif{sec:qse-nisq}) \\
$\ket{\Psi}$, $\alpha$, $\beta$ & qubit e sue ampiezze complesse (\rif{sec:qse-qubit}) \\
TVD & \ingl{total variation distance} (\rif{sec:qse-stop}) \\
$b$, $l$, $\epsilon$, $k$ & batch size, look-back window, soglia, patience di StableShots (\rif{sec:qse-stop}) \\
$\neg$, $\wedge$, $\vee$, $\rightarrow$ & negazione, \emph{e}, \emph{o}, implicazione logica \\
\bottomrule
\end{tabularx}
\end{center}
